\chapter{Introduction}\label{ch:introduction}

Complex systems are common in both nature and society and are closely related to our daily lives. Many significant large-scale phenomena, such as climate change, forest fires and pandemics, can be understood and studied as complex systems. In general, a complex system can be described as a large, coupled, nonlinear spatial system in which components interact with each other. Because of these interactions, the macroscopic behaviour of the system can be difficult to recognise from the microscopic model \cite{Fieguth2021}. 

This report will focus on self-organised criticality (SOC) and the associated non-Gaussian statistics in complex systems. In particular, I will use the Bak–Tang–Wiesenfeld (BTW) sandpile model, a representative model exhibiting self-organised criticality.

\section{Self-Organised Criticality (SOC) and Scale Invariance}

According to the research of Bak, Tang, and Wiesenfeld, dynamical systems with extended spatial degrees of freedom naturally evolve into a self-organised critical state, which is barely stable. In such a state, small perturbations at a certain location may render the neighbouring sites unstable, and the perturbations may spread to the neighbours in a chain reaction \cite{Bak1987}. Self-organised means that the system can naturally evolve to the critical state without detailed specification of the initial conditions. A critical state is a state with no characteristic time or length scales, exhibiting scale invariance \cite{Bak1988}.

Scale invariance is usually defined as the structure or property of a system remaining unchanged as the scale changes. Mathematically speaking, it can be described as $$f(\lambda x) = C(\lambda) f(x)$$
for $\lambda > 0$, where $C(\lambda)$ is the scaling factor that does not depend on $x$\cite{Khaluf2017}.

\section{Bak–Tang–Wiesenfeld (BTW) Model}

\section{Aim}


